% Spectres infrarouges de trois isomères de formule C3H6O (exercice 9, chapitre 2) ; spectres schématiques
\begin{tikzpicture}
\begin{axis}[otIR,name=a,width=75mm,height=38mm,title={Spectre A},xlabel={}]
\addplot[otSpectre] {100 \bandeIR{3005}{20}{14} \bandeIR{2925}{20}{12} \bandeIR{1715}{18}{85} \bandeIR{1420}{15}{25}
  \bandeIR{1365}{15}{50} \bandeIR{1220}{18}{60} \bandeIR{530}{15}{30}};
\end{axis}
\begin{axis}[otIR,name=b,at={(a.east)},anchor=west,xshift=12mm,width=75mm,height=38mm,title={Spectre B},xlabel={},ylabel={}]
\addplot[otSpectre] {100 \bandeIR{3330}{160}{55} \bandeIR{3085}{15}{18} \bandeIR{2900}{30}{25} \bandeIR{1645}{15}{30}
  \bandeIR{1420}{15}{20} \bandeIR{1030}{25}{70} \bandeIR{990}{12}{40} \bandeIR{925}{12}{45}};
\end{axis}
\begin{axis}[otIR,name=c,at={(a.south)},anchor=north west,yshift=-9mm,xshift=43mm,width=75mm,height=38mm,title={Spectre C}]
\addplot[otSpectre] {100 \bandeIR{2980}{25}{40} \bandeIR{2820}{12}{25} \bandeIR{2720}{15}{30} \bandeIR{1730}{18}{85}
  \bandeIR{1460}{15}{20} \bandeIR{1390}{15}{20} \bandeIR{1100}{20}{25}};
\end{axis}
\end{tikzpicture}
\par\smallskip\legende{Spectres schématiques.}
